% Assignment 3 - Arrays, braces and tables
% Tarun Konde (2025UCP1811)

\documentclass[11pt,a4paper]{article}

\usepackage[T1]{fontenc}
\usepackage{lmodern}

\usepackage[margin=2cm]{geometry}
\usepackage{amsmath}
\usepackage{amssymb}
\usepackage{array}
\usepackage{multirow}
\usepackage[table]{xcolor}
\usepackage{colortbl}
\usepackage{booktabs}
\usepackage{caption}

% colours for Que 9 and Que 10
\definecolor{headblue}{RGB}{31,78,121}
\definecolor{lightblue}{RGB}{221,235,247}
\definecolor{good}{RGB}{146,208,80}
\definecolor{avg}{RGB}{255,230,110}
\definecolor{poor}{RGB}{240,110,110}

% small helpers so the colour rule of Que 10 is written once
\newcommand{\hi}[1]{\cellcolor{good}#1}
\newcommand{\md}[1]{\cellcolor{avg}#1}
\newcommand{\lo}[1]{\cellcolor{poor}#1}

\setlength{\parindent}{0pt}

\title{Assignment 3}
\author{Tarun Konde (2025UCP1811)}
\date{\today}

\begin{document}
\maketitle

%------------------------------------------------------------------
\section*{Que 1}
\[
    \overbrace{a + b + c}^{\text{First Group}}
    + \underbrace{x + y + z}_{\text{Second Group}}
    =
    \overbrace{\underbrace{a + b}_{\text{Sum 1}} + c}^{\text{Total}}
\]

%------------------------------------------------------------------
\section*{Que 2}
\[
A =
\left[
\begin{array}{c|ccc}
         & x_1    & x_2    & x_3    \\ \hline
    y_1  & a_{11} & a_{12} & a_{13} \\
    y_2  & a_{21} & a_{22} & a_{23} \\
    y_3  & a_{31} & a_{32} & a_{33}
\end{array}
\right]
\]

%------------------------------------------------------------------
\section*{Que 3}
\[
f(x) =
\begin{cases}
    x^2 + 2x + 1,          & x < 0, \\
    \sqrt{x} + 1,          & 0 \le x < 4, \\
    \dfrac{x^2 - 1}{x - 1}, & x \ge 4.
\end{cases}
\]

%------------------------------------------------------------------
\section*{Que 4}
\[
\left[
\begin{array}{ccccc}
    \multicolumn{3}{c}{\overbrace{\hphantom{a_{11}\quad a_{12}\quad a_{13}}}^{\text{Column Group 1}}} &
    \multicolumn{2}{c}{\overbrace{\hphantom{a_{14}\quad a_{15}}}^{\text{Column Group 2}}} \\
    a_{11} & a_{12} & a_{13} & a_{14} & a_{15} \\
    a_{21} & a_{22} & a_{23} & a_{24} & a_{25} \\
    a_{31} & a_{32} & a_{33} & a_{34} & a_{35}
\end{array}
\right]
\]

%------------------------------------------------------------------
\section*{Que 5}
\[
F(x) =
\underbrace{
    \overbrace{\sum_{i=1}^{n}\left(a_i x_i^2 + b_i x_i + c_i\right)}^{\text{Quadratic Contribution}}
    +
    \overbrace{
        \begin{bmatrix} x_1 & x_2 & \cdots & x_n \end{bmatrix}
        A
        \begin{bmatrix} x_1 \\ x_2 \\ \vdots \\ x_n \end{bmatrix}
    }^{\text{Matrix Contribution}}
}_{\text{Total Energy}}
\]
where
\[
A =
\begin{bmatrix}
    a_{11} & a_{12} & \cdots & a_{1n} \\
    a_{21} & a_{22} & \cdots & a_{2n} \\
    \vdots & \vdots & \ddots & \vdots \\
    a_{n1} & a_{n2} & \cdots & a_{nn}
\end{bmatrix}.
\]

%------------------------------------------------------------------
\section*{Que 6}
\begin{center}
\begin{tabular}{|c|l|l|c|c|}
    \hline
    \textbf{Roll No.} & \textbf{Name} & \textbf{Department} & \textbf{Semester} & \textbf{CGPA} \\
    \hline
    101 & Amit  & CSE & 2 & 8.45 \\ \hline
    102 & Neha  & CSE & 2 & 9.12 \\ \hline
    103 & Rahul & ECE & 4 & 8.76 \\ \hline
    104 & Priya & ECE & 4 & 9.01 \\ \hline
\end{tabular}
\end{center}

%------------------------------------------------------------------
\section*{Que 7}
\begin{center}
\begin{tabular}{|c|c|c|}
    \hline
    \multicolumn{3}{|c|}{\textbf{Academic Performance}} \\
    \hline
    \textbf{Student} & \textbf{Midterm} & \textbf{Final} \\
    \hline
    A & 78 & 89 \\
    B & 85 & 92 \\
    C & 91 & 94 \\
    \hline
\end{tabular}
\end{center}

%------------------------------------------------------------------
\section*{Que 8}
\begin{center}
\begin{tabular}{|c|l|c|c|}
    \hline
    \textbf{Department} & \textbf{Course} & \textbf{Credits} & \textbf{Grade} \\
    \hline
    \multirow{3}{*}{CSE} & Data Structures   & 4 & A  \\ \cline{2-4}
                         & Machine Learning  & 4 & A+ \\ \cline{2-4}
                         & Operating Systems & 3 & B+ \\ \hline
    \multirow{2}{*}{ECE} & Digital Electronics & 4 & A  \\ \cline{2-4}
                         & Signals             & 3 & B+ \\ \hline
\end{tabular}
\end{center}

%------------------------------------------------------------------
\section*{Que 9}
\begin{center}
\begin{tabular}{|l|c|c|c|c|}
    \hline
    \rowcolor{headblue}
    \textcolor{white}{\textbf{Model}} &
    \textcolor{white}{\textbf{Parameters}} &
    \textcolor{white}{\textbf{FLOPs}} &
    \textcolor{white}{\textbf{Accuracy}} &
    \textcolor{white}{\textbf{Energy (kWh)}} \\
    \hline
    ResNet18        & 11.7M & 1.8G  & 94.2\% & 0.00093 \\
    \rowcolor{lightblue}
    MobileNetV2     & 3.5M  & 0.3G  & 92.8\% & 0.00061 \\
    EfficientNet-B0 & 5.3M  & 0.39G & 93.5\% & 0.00072 \\
    \hline
\end{tabular}
\end{center}

%------------------------------------------------------------------
\section*{Que 10}
Colour rule: marks $\ge 90$ green, 75--89 yellow, below 75 red.
\begin{center}
\begin{tabular}{|l|c|c|c|}
    \hline
    \textbf{Student} & \textbf{Mathematics} & \textbf{Physics} & \textbf{Computer Science} \\
    \hline
    Amit  & \hi{92}  & \md{81} & \hi{95}  \\ \hline
    Neha  & \lo{72}  & \md{88} & \hi{91}  \\ \hline
    Rahul & \lo{68}  & \lo{74}  & \md{79} \\ \hline
    Priya & \hi{96}  & \hi{94}  & \hi{98}  \\ \hline
\end{tabular}
\end{center}

%------------------------------------------------------------------
\section*{Que 11}
Fixed width \texttt{p\{\}} columns wrap the text on their own. The widths add up
to less than \verb|\textwidth| so the table fits on the page.
\begin{center}
\renewcommand{\arraystretch}{1.3}
\begin{tabular}{|>{\raggedright\arraybackslash}p{2.6cm}
                |>{\raggedright\arraybackslash}p{4cm}
                |>{\raggedright\arraybackslash}p{3.6cm}
                |>{\raggedright\arraybackslash}p{3.6cm}|}
    \hline
    \textbf{Technique} & \textbf{Description} & \textbf{Advantages} & \textbf{Limitations} \\
    \hline
    Pruning & Removes unnecessary parameters from a neural network.
            & Reduces model size and computation.
            & May reduce accuracy. \\
    \hline
    Quantization & Represents weights using lower precision.
                 & Reduces memory and energy consumption.
                 & Can introduce numerical error. \\
    \hline
    Knowledge Distillation & Transfers knowledge from a large teacher model to a smaller student model.
                           & Produces compact models.
                           & Requires training a teacher model. \\
    \hline
\end{tabular}
\end{center}

%------------------------------------------------------------------
\section*{Que 12}
Table~\ref{tab:booktabs} uses only the three \texttt{booktabs} rules and no
vertical lines.
\begin{table}[h]
    \centering
    \caption{Model comparison}
    \label{tab:booktabs}
    \begin{tabular}{lrrrr}
        \toprule
        Model & Parameters & FLOPs & Accuracy & Energy \\
        \midrule
        ResNet18        & 11.7M & 1.82G & 94.21 & 0.00093 \\
        MobileNetV2     & 3.5M  & 0.30G & 92.84 & 0.00061 \\
        EfficientNet-B0 & 5.3M  & 0.39G & 93.52 & 0.00072 \\
        \bottomrule
    \end{tabular}
\end{table}

%------------------------------------------------------------------
\section*{Que 13}
% The inner tables only have the lines between their own cells, so they
% do not double up with the borders of the outer table. \inner adds a
% little space above and below each one.
\newcommand{\inner}[2]{%
    \begin{tabular}[c]{c|c}
        \rule{0pt}{2.6ex}Metric & Value \\ \hline
        Accuracy & #1 \\
        F1 Score & #2
    \end{tabular}}
\begin{center}
\renewcommand{\arraystretch}{1.2}
\begin{tabular}{|c|c|c|}
    \hline
    \textbf{Model} & \textbf{Architecture} & \textbf{Performance} \\
    \hline
    CNN    & Conv $\rightarrow$ ReLU $\rightarrow$ Pool & \inner{94\%}{93\%} \\[0.6ex] \hline
    ResNet & Conv $\rightarrow$ BN $\rightarrow$ ReLU   & \inner{96\%}{95\%} \\[0.6ex] \hline
    ViT    & Attention $\rightarrow$ MLP                 & \inner{97\%}{96\%} \\[0.6ex] \hline
\end{tabular}
\end{center}

%------------------------------------------------------------------
\section*{Que 14}
\begin{center}
\renewcommand{\arraystretch}{2.2}
\begin{tabular}{|l|c|c|c|}
    \hline
    \textbf{Model} & \textbf{Loss Function} & \textbf{Parameters} & \textbf{Complexity} \\
    \hline
    Linear      & $\displaystyle L = \frac{1}{n}\sum_{i=1}^{n}(y_i - \hat{y}_i)^2$ & 10K  & $O(n)$ \\
    \hline
    CNN         & $\displaystyle L = -\sum y \log(\hat{y})$                         & 1.2M & $O(n^2)$ \\
    \hline
    Transformer & $\displaystyle L = -\sum_i y_i \log p_i$                          & 110M & $O(n^2 d)$ \\
    \hline
\end{tabular}
\end{center}

\end{document}
